% !TEX root = ../notes.tex

\documentclass[../notes.tex]{subfiles}

\begin{document}

\section{March 30}

Welcome back.

\subsection{Fun with \texorpdfstring{$K(1)$}{ K(1)}}
We are currently talking about Bousfield localizations. To get started, let's discuss $\mathrm{BP}$-localization.
\begin{remark}
	Note that \Cref{prop:bp-bousfield} implies that $L_{\mathrm{BP}}=L_{(K(0)\oplus K(1)\oplus\cdots)\oplus\FF_p}$.
\end{remark}
\begin{remark}
	If $X$ is bounded below, then $L_{\mathrm{BP}}X=X$. Indeed, because $X$ is bounded below, one finds that the Adams--Novikov spectral sequence converges, which means
	\[X=\lim_\Delta\left(X\otimes\mathrm{BP}^{\otimes(\bullet+1)}\right).\]
	Thus, $X$ is a limit of $\mathrm{BP}$-modules, and $L_{\mathrm{BP}}$ does nothing to such modules.
\end{remark}
In fact, one can often recover a module from its localizations, as we saw in \Cref{thm:recover-x-from-localizations}.
% \begin{theorem}[Hopkins--Ravenel]
% 	If $X$ is a perfect $\mathbb S_{(p)}$-module, then
% 	\[X=\lim_nL_nX.\]
% \end{theorem}
\begin{example}
	On the homework, we will show that $\lim_nL_n(\ZZ)\ne\ZZ$. In fact, $\lim_nL_n(\ZZ)=\QQ$.
\end{example}
\begin{remark}
	On the homework, we will show that $L_{\mathrm{BP}}L_n^f=L_n$.
\end{remark}
Last time, we showed that $L_n$ is smashing (see \Cref{thm:ln-smashing}), assuming the fact that $L_{K(n)}$ of a type-$n$ complex is in the thick subcategory generated by $K(n)$ (written in \Cref{thm:galois-theory}).

Our next goal is to explain how to compute $L_{K(n)}X$, for example in attempt to prove \Cref{thm:galois-theory}. We will also eventually want to approach \Cref{conj:chromatic-splitting-again}.
\begin{example}
	Note that $L_{K(0)}X=L_0X$, which is $L_0\mathbb S\otimes X$ because $L_0$ is smashing. Now, $L_0\mathbb S=\QQ$, so this is $\QQ\otimes X$, which is not too hard to understand (e.g., by homotopy groups) because $\QQ$ is a field.
\end{example}
We now move on to $K(1)$. This turns out to be $\mathrm{KU}/p$, where we recall that $\mathrm{KU}=\Sigma_+^\infty\CP^\infty[1/\beta]$ so that $\pi_*\mathrm{KU}=\ZZ_{(p)}[\beta,1/\beta]$.
\begin{example}
	Take $p=2$. Then we compute $L_{K(1)}\Sigma^\infty_+\RP^\infty$. To start, we should compute $\mathrm{KU}^*\RP^\infty$. Well, there is a fiber sequence $\RP^\infty\to\CP^\infty\to\CP^\infty$, where the second map is given by $\mc L\mapsto\mc L^{\otimes2}$. In particular, it follows that $\RP^\infty=\colim_{\CP^\infty}\CP^\infty$. Now, $\mathrm{KU}^*\CP^\infty=\pi_*\mathrm{KU}[[x]]$ (with formal group law structure), so the Atiyah--Hirzebruch specral sequence tells us that
	\[\mathrm{KU}^*\RP^\infty\cong\frac{\ZZ_{(2)}[\beta,1/\beta][[x]]}{[2](x)}.\]
	Now, the formal group law has $[2](x)=2x+x^2$. Taking the quotient by $p$ then gives
	\[K(1)^*\RP^\infty\cong\frac{\FF_2[\beta,1/\beta][x]}{\left(x^2\right)}.\]
	Thus, we see that $K(1)\otimes\Sigma^\infty_+\RP^\infty=K(1)\oplus K(1)$ (by computing homotopy groups and using the fact that $K(1)$ is a skew field), so $K(1)\otimes\Sigma^\infty\RP^\infty=K(1)$.
\end{example}
\begin{remark}
	Recall that there is a transfer map $\Sigma_+^\infty\RP^\infty\to\mathbb S$. It turns out that the induced map $\Sigma^\infty\RP^\infty\to\mathbb S$ is an isomorphism in $K(1)$-homology, so it follows that $L_{K(1)}\Sigma^\infty\RP^\infty=L_{K(1)}\mathbb S$. Thus, we have computed $L_{K(1)}\mathbb S$.
\end{remark}
\begin{remark}
	It is a theorem of Ravenel that there is an idempotent $z$ for which $L_{K(1)}\Sigma_+^\infty\RP^\infty[1/z]=L_{K(1)}\mathbb S$.
\end{remark}
While we're here, let's discuss the picture for odd primes $p$.
\begin{example}
	At odd primes $p$, it turns out that $L_{K(1)}\Sigma_+^\infty\mathrm BC_p=L_{K(1)}^{\oplus p}$, and there is an idempotent $z$ for which $L_{K(1)}\Sigma_+^\infty\mathrm BC_p[1/z]=L_{K(1)}\mathbb S^{\oplus(p-1)}$. Here, $\mathrm BC_p=K(\ZZ/p\ZZ,1)$, so it admits an action by $\FF_p^\times$. It follows that we can recover $L_{K(1)}\mathbb S$ by taking invariants: one has
	\[L_{K(1)}\mathbb S=\left(L_{K(1)}\Sigma_+^\infty\mathrm BC_p[1/z]\right)^{h\FF_p^\times},\]
	where $(-)^{h\FF_p^\times}$ means that we are taking the limit from $\mathrm B\FF_p^\times$.
\end{example}
\begin{example} \label{ex:get-a-galois-extension}
	We continue the previous example. Note that any $k>0$ admits an embedding $\ZZ/p\ZZ\to\ZZ/p^k\ZZ$, so there is an $\mathbb E_\infty$-ring map $\Sigma_+^\infty\mathrm B(\ZZ/p\ZZ)\to\Sigma_+^\infty\mathrm B(\ZZ/p^k\ZZ)$. Taking $L_{K(1)}$, we receive a map
	\[L_{K(1)}\Sigma^\infty_+\mathrm B(\ZZ/p\ZZ)[1/z]\to L_{K(1)}\Sigma_+^\infty\left(\ZZ/p^k\ZZ\right)[1/z]\]
	of $\mathbb E_\infty$-rings. It is a theorem that
	\[L_{K(1)}\mathbb S=\left(L_{K(1)}\Sigma_+^\infty\mathrm B\left(\ZZ/p^k\ZZ\right)\right)^{h\left(\ZZ/p^k\ZZ\right)^\times}.\]
	(Namely, one can compute that this is an isomorphism on $K(1)$-homology, so it becomes an isomorphism after taking $L_{K(1)}$.)
\end{example}

\subsection{Galois Extensions}
These sorts of maps to fixed points are interesting enough to be worth a name.
\begin{definition}[Galois extension]
	Fix a symmetric monoidal stable category $\mc C$. Then an $\mathbb E_\infty$-algebra $A$ is a \textit{faithful Galois} extension of the unit if and only if there is a finite group $G$ acting on $A$ via $\mathbb E_\infty$-automorphisms, satisfying the following.
	\begin{itemize}
		\item The natural map $1\to A^{hG}$ is an isomorphism.
		\item The natural map $A\otimes A\to\prod_GA$ given by $a_1\otimes a_2\mapsto(a_1g(a_2))_{g\in G}$ is an isomorphism.
		\item An element $X\in\mc C$ is zero if and only if $X\otimes A$ is zero.
	\end{itemize}
\end{definition}
\begin{remark}
	Motivated by \Cref{ex:get-a-galois-extension}, it is a theorem that $L_{K(1)}\Sigma_+^\infty\mathrm B\left(\ZZ/p^k\ZZ\right)$ is a Galois extension of $L_{K(1)}\mathbb S$ in the category of $K(1)$-local spectra. (The tensor product in $K(1)$-local spectra must localize by $K(1)$ afterward.)
\end{remark}
The moral is that we have an infinite ``chain'' of Galois extensions of $L_{K(1)}\Sigma_+^\infty\mathrm B\left(\ZZ/p^k\ZZ\right)$ over $L_{K(1)}\mathbb S$. We can even try to take a union to produce an infinite Galois extension: namely, we can take the colimit
\[\colim_k\Sigma^\infty_+\mathrm B\left(\ZZ/p^k\ZZ\right)=\Sigma_+^\infty\mathrm B^2\ZZ_p\]
up to $p$-completion, which one sees from a direct calculation at finite level. But this target is just the completion of $\Sigma^\infty_+\mathrm B^2\ZZ=\Sigma^\infty_+\CP^\infty$. Now, taking $L_{K(1)}$, we see we are looking at $L_{K(1)}\Sigma_+^\infty\mathrm B^2\ZZ_p=L_{K(1)}\Sigma_+^\infty\CP^\infty$, which more manifestly has a $\ZZ_p^\times$-action.

Accordingly, we would like to hunt for a $\ZZ_p^\times$-action on $L_{K(1)}\Sigma^\infty_+\CP^\infty[1/z]$ so that we receive a Galois extension of $L_{K(1)}\mathbb S$.
\begin{remark}
	In $\pi_*L_{K(1)}\Sigma^\infty_+\CP^\infty$, it turns out that $z$ is a unit multiplied by $\Sigma^\infty\beta$.
\end{remark}
Thus, by passing the localization through, we see that the $\ZZ_p^\times$-Galois extension should be the $K(1)$-local\-ization of $\Sigma_+^\infty\CP^\infty[1/\beta]$, which is just $\mathrm{KU}$.
\begin{remark}
	We claim that $L_{K(1)}\mathrm{KU}=\mathrm{KU}^\land_p$. Indeed, $L_{K(1)}\mathrm{KU}$ needs to at least be $p$-complete, so it only remains to check that $\mathrm{KU}^\land_p$ is actually $K(1)$-local. But this is true because $\mathrm{KU}^\land_p$ is the limit of the modules $\mathrm{KU}/p^n$, and $K(1)=\mathrm{KU}/p$.
\end{remark}
It will shortly turn out that the action of $\ZZ_p^\times$ on $\mathrm{KU}^\land_p$ will define a Galois extension of $L_{K(1)}\mathbb S$, as soon as we can make sense of what that means.
\begin{remark}
	If $\ell$ is a topological generator of $\ZZ_p^\times$ (which exists for $p>2$), then we find that taking $\ZZ_p^\times$-invariants amounts to taking the fiber of $\psi^\ell\colon\mathrm{KU}^\land_p\to\mathrm{KU}^\land_p$. In particular, we see that
	\[L_{K(1)}\mathbb S=\op{fib}\left(\psi^\ell\colon\mathrm{KU}^\land_p\to\mathrm{KU}^\land_p\right).\]
\end{remark}
The previous remark implies (by modding out by $p$) that $L_{K(1)}\mathbb S/p$ is the fiber of the map $K(1)\to K(1)$. Thus, we have proven a special case of \Cref{thm:galois-theory}.

In general, one hopes that $L_{K(n)}\mathbb S$ looks like $E^{h\mathbb G}$, where $\mathbb G$ is some profinite group acting continuously on some even $\mathbb E_\infty$-algebra related to $K(n)$.
\begin{remark}
	Let's explain what happens for higher heights $n$. Here, $L_{K(n)}\mathbb S$ turns out to admit the Galois extensions $L_{K(n)}K\left(\ZZ/p^k\ZZ,n\right)[1/z_n]$ with Galois group $\left(\ZZ/p^k\ZZ\right)^\times$, which we can assemble into a Galois extension $L_{K(n)}K(\ZZ_p,n+1)[1/z_n]$ with Galois group $\ZZ_p^\times$. However, this cannot be the end of the story because these $\mathbb E_\infty$-algebras are not even!
\end{remark}
Thus, we need to find some higher Galois extensions. These will come from a functor $E$ which takes pairs $(R_0,\Gamma_0)$, where $R_0$ is a (classical) perfect $\FF_p$-algebra and a formal group $\Gamma_0$ over $R_0$, and it outputs an even $\mathbb E_\infty$-ring $E(R_0,\Gamma_0)$ whose homotopy ring in degree zero is a ``deformation'' of $R_0$. Furthermore, the formal group attached to $E(R_0,\Gamma_0)$ (as an even $\mathbb E_\infty$-ring) will roughly be a deformation of $\Gamma_0$.
\begin{example}
	It will turn out that $E(\FF_p,\mathbb G_m)=\mathrm{KU}^\land_p$, whose $\pi_0$ is $\ZZ_p$.
\end{example}
\begin{remark}
	Let's explain the application to the chromatic splitting conjecture. We will be able to compute with $L_{K(n-1)}L_{K(n)}\mathbb S^{\mathrm{cyc}}$, where $L_{K(n)}\mathbb S_p^{\mathrm{cyc}}$ is $L_{K(n)}\Sigma_+^\infty K(\ZZ_p,n+1)[1/z]$ is a $\ZZ_p^\times$-extension of $L_{K(n)}\mathbb S$. In particular, the former object is a Galois extension, where we find that
	\[L_{K(n-1)}L_{K(n)}\mathbb S^{\mathrm{cyc}}=E(R_0,\Gamma_0)^{h\mathbb G},\]
	and the chromatic splitting map lifts this to a map to $E(\widehat R_0,\Gamma_0)^{h\mathbb G}$, where $\widehat R_0=\FF_p((t^{1/p^\infty}))$. Now, $\widehat R_0$ is perfect, so we are able to do number theory to it, which helps us with some cohomology calculations.
\end{remark}

\end{document}